\documentclass[../principal.tex]{subfiles}
\graphicspath{{\subfix{../imagenes/}}}

\begin{document}

\ifSubfilesClassLoaded{
    \chapter{Componentes básicos}
}{}

\section{Capacitores}

Los capacitores son componentes pasivos de dos terminales que se utilizan para almacenar energía eléctrica en un circuito.

Se representan con el símbolo $C$ y se miden en Farads, cuyo símbolo es la letra F mayúscula ($\mathrm{F}$).

En este libro usamos la palabra capacitor para denominar al componente, y la palabra capacitancia para referirnos a la propiedad física del material, entonces, el capacitor tiene capacitancia.

En español muchas veces se usa la palabra capacitancia para ambos conceptos, pero preferimos hacer la distinción, tal como se hace en inglés entre capacitor y capacitance.

También en español se tiende a usar la palabra condensador, pero nos gusta en este libro que el nombre del componente sea lo más parecido posible a su propiedad física, entonces preferimos usar la palabra capacitor.

Así los resistores tienen o ejercen resistencia, y los capacitores tienen o ejercen capacitancia.

\subsection{En serie}

Si tenemos dos capacitores en serie, llamados $C_1$ y $C_2$, entre los nodos A y B, podemos diagramarlos así:

\begin{circuitikz}
\draw
  (0,0) node[left]{A}
  to[short, o-] (1,0)
  to[C, l=$C_1$] (3,0)
  to[C, l=$C_2$] (5,0)
  to[short, -o] (6,0) node[right]{B};
\end{circuitikz}

Podemos reemplazar estos dos capacitores por uno equivalente $C_{eq}$, que se calcula usando la fórmula:

\begin{equation}
\frac{1}{C_{eq}} = \frac{1}{C_1} + \frac{1}{C_2}
\end{equation}

o usando la fórmula equivalente:

\begin{equation}
\frac{1}{C_{eq}} = \frac{1}{C_1} + \frac{1}{C_2}
\end{equation}

y diagramado así:

\begin{circuitikz}
\draw
  (0,0) node[left]{A}
  to[short, o-] (2,0)
  to[C, l=$C_{eq}$] (4,0)
  to[short, -o] (6,0) node[right]{B};
\end{circuitikz}

\subsection{En paralelo}

Si tenemos dos capacitores en paralelo, llamados $C_1$ y $C_2$, entre los nodos A y B, podemos diagramarlos así:

\begin{circuitikz}
% nodos A y B
\coordinate (A) at (0,0);
\coordinate (B) at (6,0);

% dibujar etiquetas de los nodos
\draw (A) node[left]{A};
\draw (B) node[right]{B};

% puntos de conexion
\coordinate (izq) at (1,0);
\coordinate (der) at (5,0);

% rama arriba
\draw (A) to [short, o-] (izq)
      (izq) to[short] ++(0,1.0)
      to[C, l=$C_1$] ++(4,0)
      to[short] (der)
      (der) to [short, -o] (B);

% rama abajo
\draw (A) to [short, o-] (izq)
      (izq) to[short] ++(0,-1.0)
      to[C, l=$C_2$] ++(4,0)
      to[short] (der)
      (der) to [short, -o] (B);

\end{circuitikz}

Podemos reemplazar estos dos capacitores por uno equivalente $C_{eq}$, que se calcula usando la fórmula:

\begin{equation}
C_{eq} = C_1 + C_2
\end{equation}

y diagramado así:

\begin{circuitikz}
\draw
  (0,0) node[left]{A}
  to[short, o-] (2,0)
  to[C, l=$C_{eq}$] (4,0)
  to[short, -o] (6,0) node[right]{B};
\end{circuitikz}

\newpage

\end{document}


